% 第6章 6.2.3：同一函数 f=(3,1,2) 在三组边权下的沿边变化能量。
% 边差 (2,-1)，平方 (4,1)；能量 = w12*4 + w23*1：
% (1,1)->5，(1,4)->8，(1/4,1)->2。线宽与边权成正比（示意）。
\begin{tikzpicture}[x=1mm,y=1mm,font=\figurefont\footnotesize,text=FigureInk,
  vertex/.style={circle,draw=FigureStroke,fill=FigureBlueFill,line width=.8pt,minimum size=6mm,inner sep=0pt},
  note/.style={font=\FigureNoteFont\fontsize{7.5}{9}\selectfont,text=FigureMuted}]
\foreach \dx/\wa/\wb/\la/\lb/\E/\tag in {0/1/1/1/1/5/(a),38/1/4/1/4/8/(b),76/0.25/1/{1/4}/1/2/(c)}{
 \begin{scope}[xshift=\dx mm]
  \node[vertex] (a) at (3,12) {$1$};
  \node[vertex] (b) at (17,12) {$2$};
  \node[vertex] (c) at (31,12) {$3$};
  \draw[FigureStroke,line width=\wa*0.9pt] (a)--(b);
  \draw[FigureStroke,line width=\wb*0.9pt] (b)--(c);
  \node[text=FigureBlue] at (3,19) {$3$};
  \node[text=FigureBlue] at (17,19) {$1$};
  \node[text=FigureBlue] at (31,19) {$2$};
  \node[note] at (10,15.2) {$w=\la$};
  \node[note] at (24,15.2) {$w=\lb$};
  \node at (17,3.5) {$\la\cdot 4+\lb\cdot 1=\E$};
  \node[anchor=west] at (-1.5,25) {\tag};
 \end{scope}}
\end{tikzpicture}
